\section{Discussion}
\paragraph{Shared operations, task-dependent organization.}
The atlas connects domains through operations that differ in prevalence, timing, and function: verification checks derivations or action preconditions, while revision rejects unoriginal ideas. Matched instructions connect some profiles to task demand. This offers a behavioral connection to the organization of simpler processes in \citet{minsky1986society} and flexible recruitment of shared systems \citep{cole2013hubs}. The evidence concerns textual organization, not discrete mental agents or corresponding brain regions.

\paragraph{Description does not supply a prescription.}
Domain information is robust within the annotated population, but it neither supplies supported conditional outcome prediction nor ensures faithful control. The same-domain GPQA failure is especially instructive: the vector does not reliably induce its intended contrast. Its harm therefore cannot establish that the natural profile is undesirable. Large budget-dependent changes in answer production further constrain what the transfer matrix means. Resource-rational accounts motivate evaluating benefits against computation costs \citep{shenhav2013control,lieder2020resource}; our results do not identify an optimal allocation policy.

\paragraph{Explanation requires more than a similar average shift.}
A2, the saved zero-start RL vector, and matched SEAL all produce similarly directed average shifts, with different magnitudes and outcomes. Their resemblance is therefore compatible with a common coarse description of shorter or reorganized reasoning. It does not identify distinct mechanisms arriving at the same useful computation. This also suggests a boundary question for \citet{kim2026society}: behaviors associated with stronger models or their acquisition need not be beneficial to increase within an already trained reasoner. Our bias-only interventions, measurement scale, and historical training audit preclude interpreting that question as a refutation of full-model learning results.

\paragraph{Connections across objectives.}
The math-profile vector lowers paired originality, whereas a math-style instruction raises it while reducing feasibility. Their different effects rule out treating the two as equivalent implementations of mathematical thinking. They do not establish a universal opposition between analysis and creativity; human creative cognition involves cooperation between spontaneous and controlled processes \citep{beaty2015creative}. Cross-task evaluation makes completion and judged quality distinct consequences of reasoning control.

\paragraph{Hypotheses and future work.}
The findings motivate three hypotheses. Transfer may depend on unmet computational demands rather than source-domain similarity; this calls for a selection rule frozen before fresh evaluation, with calibrated dose, answer elicitation, and stopping controls. An operation's value may depend on reasoning stage; event-triggered interventions could test error repair and candidate generation separately from termination. Finally, similar aggregate labels may hide differences in effective verification or candidate diversity. Cross-family annotation and functional coding can test whether a finer representation explains outcomes beyond duration. These hypotheses extend the connection to flexible task organization without assuming corresponding brain systems.

\section{Conclusion}
Across seven models, annotated operation frequencies and timing distinguish benchmark conditions, with some differences responsive to matched task demands. Their use as intervention targets is less secure: fidelity varies, large outcome changes accompany answer production, and similar average shifts leave important differences unresolved. Thought Atlas makes this separation between description, prescription, and explanation explicit.
